\documentclass{article}
\usepackage{amsmath, amssymb, amsthm}
\usepackage{hyperref}
\usepackage{geometry}
\geometry{margin=1in}

\title{Erd\H{o}s Problem \#1053}
\date{Last updated: 2025-09-28}
\author{Source: erdosproblems.com}

\begin{document}
\maketitle

\noindent\textbf{Status:} OPEN \quad \textbf{No prize}

\noindent\textbf{Tags:} number theory

\medskip
\noindent\textbf{Problem Statement:}

Call a number $k$-perfect if $\sigma(n)=kn$, where $\sigma(n)$ is the sum of the divisors of $n$. Must $k=o(\log\log n)$?

\bigskip
\noindent\textbf{Remarks:}

A question of Erd\H{o}s, as reported in problem B2 of Guy's collection \cite{Gu04}. Guy further writes 'It has even been suggested that there may be only finitely many $k$-perfect numbers with $k\geq 3$.' The largest $k$ for which a $k$-perfect number has been found is $k=11$ - see this page for more information.

These are known as multiply perfect numbers. When $k=2$ this is the definition of a perfect number.

\bigskip
\noindent\textbf{References:}

[Gu04] Guy, Richard K., Unsolved problems in number theory. (2004), xviii+437.

\bigskip
\noindent\small{Source: \url{https://www.erdosproblems.com/1053}}
\end{document}
